\documentclass[notes=show]{beamer}
%%%%%%%%%%%%%%%%%%%%%%%%%%%%%%%%%%%%%%%%%%%%%%%%%%%%%%%%%%%%%%%%%%%%%%%%%%%%%%%%%%%%%%%%%%%%%%%%%%%%%%%%%%%%%%%%%%%%%%%%%%%%%%%%%%%%%%%%%%%%%%%%%%%%%%%%%%%%%%%%%%%%%%%%%%%%%%%%%%%%%%%%%%%%%%%%%%%%%%%%%%%%%%%%%%%%%%%%%%%%%%%%%%%%%%%%%%%%%%%%%%%%%%%%%%%%
\AtBeginSubsection[
  {\frame<beamer>{\frametitle{Outline}   
    \tableofcontents[currentsection,currentsubsection]}}%
]%
{
  \frame<beamer>{ 
    \frametitle{Outline}   
    \tableofcontents[currentsection,currentsubsection]}
}
\usepackage{ulem}
\usepackage{pdflscape}
\setbeamertemplate{caption}{\insertcaption}
\usepackage{color,xcolor,ucs}
\usepackage{beamerthemesplit}
\usepackage{graphicx}
\usepackage{amsmath}
\usepackage{adjustbox}
\usepackage{pdfpages} 
\definecolor{mintgreen}{rgb}{0.6, 1.0, 0.6}
\makeatletter
\@addtoreset{subfigure}{framenumber}% subfigure counter resets every frame
\makeatother
\newcommand*{\Scale}[2][4]{\scalebox{#1}{$#2$}}%
\newcommand*{\Resize}[2]{\resizebox{#1}{!}{$#2$}}%
\usepackage{amsfonts}
\usepackage{booktabs}
\usepackage{color,soul}
\usepackage{pdfpages}
\setboolean{@twoside}{false}
\usepackage{amsmath}
\usepackage{pdflscape} 
\usepackage{eurosym}
\usepackage{amstext} % for \text
\DeclareRobustCommand{\officialeuro}{%
  \ifmmode\expandafter\text\fi
  {\fontencoding{U}\fontfamily{eurosym}\selectfont e}}
\usepackage[caption = false]{subfig}
\usepackage{appendixnumberbeamer} 
\usepackage{graphicx}
\usepackage{mathpazo}
\usepackage{hyperref}
\usepackage{multimedia}
\usepackage{graphicx}
\usepackage{multirow}
\usepackage{subfig}
\usepackage{verbatim}
\usepackage{booktabs, multicol, multirow}
\setcounter{MaxMatrixCols}{10}
\input{Macros.tex}
\AtBeginSection[]
{
  \begin{frame}[noframenumbering]
    \frametitle{Outline for section \thesection}
    \tableofcontents[currentsection]
  \end{frame}
}


\usepackage{xcolor,soul}
\renewcommand<>{\hl}[1]{\only#2{\beameroriginal{\hl}}{#1}}

% https://tex.stackexchange.com/questions/41683/why-is-it-that-coloring-in-soul-in-beamer-is-not-visible
\makeatletter
\newcommand\SoulColor{%
  \let\set@color\beamerorig@set@color
  \let\reset@color\beamerorig@reset@color}
\makeatother
\SoulColor
\AtBeginSection{\frame{\sectionpage}}

\newsavebox\mybox
\newenvironment{aquote}[1]
  {\savebox\mybox{#1}\begin{quote}\openautoquote\hspace*{-.7ex}}
  {\unskip\closeautoquote\vspace*{1mm}\signed{\usebox\mybox}\end{quote}}

\newtheorem{proposition}[theorem]{Proposition}
\newtheorem{Assumption}[theorem]{Assumption}
\newenvironment{stepenumerate}{\begin{enumerate}[<+->]}{\end{enumerate}}
\newenvironment{stepitemize}{\begin{itemize}[<+->]}{\end{itemize} }
\newenvironment{stepenumeratewithalert}{\begin{enumerate}[<+-| alert@+>]}{\end{enumerate}}
\newenvironment{stepitemizewithalert}{\begin{itemize}[<+-| alert@+>]}{\end{itemize} }
\usetheme{Madrid}
\makeatletter
\setbeamertemplate{footline}
{
  \leavevmode%
  \hbox{%
  \begin{beamercolorbox}[wd=.333333\paperwidth,ht=2.25ex,dp=1ex,center]{author in head/foot}%
    \usebeamerfont{author in head/foot}\insertshortauthor~~\beamer@ifempty{\insertshortinstitute}{}{(\insertshortinstitute)}
  \end{beamercolorbox}%
  \begin{beamercolorbox}[wd=.333333\paperwidth,ht=2.25ex,dp=1ex,center]{title in head/foot}%
    \usebeamerfont{title in head/foot}\insertshorttitle
  \end{beamercolorbox}%
  \begin{beamercolorbox}[wd=.333333\paperwidth,ht=2.25ex,dp=1ex,right]{date in head/foot}%
    \usebeamerfont{date in head/foot}\insertshortdate{}\hspace*{2em}
%    \insertframenumber{} / \inserttotalframenumber\hspace*{2ex} % DELETED
  \end{beamercolorbox}}%
  \vskip0pt%
}
\makeatother
\usepackage{pdfpages}
\begin{document}
	\title[]{\large Instrumental Variables
	\vspace{0.5cm}
	}
	\author[Raffaele Saggio - Econ 628]{Raffaele Saggio}
	\institute[]{UBC}
	\date{\today \\
}
	\maketitle






\begin{frame}{Introduction}

\begin{itemize}
\item This lecture covers \textbf{instrumental variables} (IV) methods{\small{}\medskip{}
}{\small\par}
\item IV has been at the center of applied microeconometric research for
the last few decades{\small{}\medskip{}
}{\small\par}
\item IV methods pioneered by economists have been widely adopted in other
areas of empirical research{\small{}\medskip{}
}{\small\par}
\item We will begin by studying IV methods in a constant effects framework,
then circle back to reinterpret these methods in a world of heterogeneous
effects
\end{itemize}
\end{frame}

\begin{frame}{IV History}

\begin{itemize}
\item {\small{}Instrumental variables methods were developed to deal with
two seemingly unrelated problems\medskip{}
}{\small\par}
\item {\small{}Philip and Sewall Wright, 1920s: \medskip{}
}

\begin{itemize}
\item {\small{}Sought to estimate supply and demand relationships with data
on equilibrium prices/quantities\medskip{}
}{\small\par}
\item {\small{}This line of inquiry developed into the study of simultaneous
equations models (SEMs)\medskip{}
}{\small\par}
\end{itemize}
\item {\small{}Wald, 1940: Sought to estimate a regression line when the
regressor is measured with error\medskip{}
}{\small\par}
\item {\small{}Today IV is most often used as a method to eliminate omitted
variables bias in single-equation causal models}{\small\par}
\end{itemize}
\end{frame}
\begin{frame}
\frametitle{Outline}
\begin{enumerate}
\item Part 1: Constant Treatment Effect
\begin{itemize}
\item Simple IV: 1 instrument, 1 endogenous variable
\item Multiple Instruments: 2SLS 
\item IV as GMM
\item Overidentification Tests
\item 2SLS Mistakes
\item Pop-Quiz!
\item Weak IV (Appendix)
\end{itemize}
\bigskip
\item Part 2: Heterogenous Treatment Effects
\begin{itemize}
\item LATE Theorem
\item Compliers Characteristics 
\item Over-id tests with HTE
\end{itemize}
\bigskip
\item Part 3: IV and Sample Selection Models
\begin{itemize}
\item Heckit 
\item Control Function
\item Lee Bounds
\end{itemize}
\end{enumerate}
\bigskip
Main reference: Mostly Harmless Econometrics. 
\end{frame}

\section{Constant Effects}
\subsection{Simple IV}
\begin{frame}{A Simple Causal Model}

\begin{itemize}
\item {\small{}Consider a simple constant-effects causal model:}{\small\par}
\end{itemize}
\begin{center}
{\small{}$Y_{i}(s)=\alpha+\rho s+\epsilon_{i}$}{\small\par}
\par\end{center}
\begin{itemize}
\item {\small{}$Y_{i}(s)$ is $i$'s potential wage if s/he attends school
for $s$ years\medskip{}
}{\small\par}
\item {\small{}This additively separable function form implies constant
causal effects: $Y_{i}(s)-Y_{i}(s-1)=\rho\ \forall(i,s)$\medskip{}
}{\small\par}
\item {\small{}We observe schooling $S_{i}$ and earnings $Y_{i}(S_{i})$}{\small\par}
\end{itemize}
\end{frame}
\begin{frame}{A Simple Causal Model}

\begin{itemize}
\item {\small{}Equation for observed earnings:}{\small\par}
\end{itemize}
\begin{center}
{\small{}$Y_{i}=\alpha+\rho S_{i}+\epsilon_{i}$}{\small\par}
\par\end{center}
\begin{itemize}
\item {\small{}This is a structural equation, not a regression, so no guarantee
that $Cov(S_{i},\epsilon_{i})=0$\medskip{}
}{\small\par}
\item {\small{}Project $\epsilon_{i}$ on to unobserved ability $A_{i}$:}{\small\par}
\end{itemize}
\begin{center}
{\small{}$\epsilon_{i}=\gamma A_{i}+v_{i}$}{\small\par}
\par\end{center}
\begin{itemize}
\item {\small{}$Cov(A_{i},v_{i})=0$ by definition; assume $Cov(S_{i},v_{i})=0$\medskip{}
}{\small\par}
\item {\small{}Unobserved ability $A_{i}$ is a potential source of omitted
variable bias}{\small\par}
\end{itemize}
\end{frame}


\begin{frame}{Long and Short Regressions}

\begin{itemize}
\item {\small{}``Long'' regression:}{\small\par}
\end{itemize}
\begin{center}
{\small{}$Y_{i}=\alpha+\rho S_{i}+\gamma A_{i}+v_{i}$}{\small\par}
\par\end{center}
\begin{itemize}
\item {\small{}``Short'' regression:}{\small\par}
\end{itemize}
\begin{center}
{\small{}$Y_{i}=\alpha_{s}+\rho_{s}S_{i}+v_{is}$}{\small\par}
\par\end{center}
\begin{itemize}
\item {\small{}OVB formula:}{\small\par}
\end{itemize}
\begin{center}
{\small{}$\rho_{s}=\rho+\gamma\dfrac{Cov(S_{i},A_{i})}{Var(S_{i})}$}{\small\par}
\par\end{center}
\begin{itemize}
\item {\small{}The regression we've got (short) is not the regression we
want (long)\medskip{}
}{\small\par}
\item {\small{}IV is a method for obtaining the long regression coefficient
without data on the omitted regressor $A_{i}$}{\small\par}
\end{itemize}
\end{frame}

\begin{frame}{Instrumental Variables}

\begin{center}
{\small{}$Y_{i}=\alpha+\rho S_{i}+\epsilon_{i}$}{\small\par}
\par\end{center}
\begin{itemize}
\item {\small{}Suppose we find an }\textbf{\small{}instrument}{\small{},
$Z_{i}$, that satisfies two conditions:\medskip{}
}

\begin{enumerate}
\item \textbf{\small{}First stage: }{\small{}$Cov(S_{i},Z_{i})\neq0$\medskip{}
}{\small\par}
\item \textbf{\small{}Exclusion restriction: }{\small{}$Cov(\epsilon_{i},Z_{i})=0$\medskip{}
}{\small\par}
\end{enumerate}
\item {\small{}The first stage condition requires the instrument to be correlated
with the causal variable of interest. This is testable.}{\small\par}
\end{itemize}
\end{frame}

\begin{frame}{Instrumental Variables}

\begin{center}
{\small{}$Y_{i}=\alpha+\rho S_{i}+\epsilon_{i}$}{\small\par}
\par\end{center}
\begin{enumerate}
\item \textbf{\small{}First stage: }{\small{}$Cov(S_{i},Z_{i})\neq0$\medskip{}
}{\small\par}
\item \textbf{\small{}Exclusion restriction: }{\small{}$Cov(\epsilon_{i},Z_{i})=0$\medskip{}
}{\small\par}
\end{enumerate}
\begin{itemize}
\item {\small{}The exclusion restriction requires the instrument to be uncorrelated
with the structural error $\epsilon_{i}$. This requires\medskip{}
}

\begin{itemize}
\item {\small{}$Z_{i}$ must be uncorrelated with all unobserved determinants
of earnings (``as good as randomly assigned'')\medskip{}
}{\small\par}
\item {\small{}$Z_{i}$ has no direct causal impact on $Y_{i}$\medskip{}
}{\small\par}
\end{itemize}
\item {\small{}The exclusion restriction is not testable (caveats later)}{\small\par}
\end{itemize}
\end{frame}

\begin{frame}{Instrumental Variables}

\begin{itemize}
\item {\small{}The exclusion restriction implies}{\small\par}
\end{itemize}
\begin{center}
{\small{}$Cov\left(Y_{i},Z_{i}\right)=Cov\left(\alpha+\rho S_{i}+\epsilon_{i},Z_{i}\right)$}{\small\par}
\par\end{center}

\begin{center}
{\small{}$=\rho Cov(S_{i},Z_{i})$}{\small\par}
\par\end{center}
\begin{itemize}
\item {\small{}If the first stage assumption is satisfied, we can divide
by $Cov(S_{i},Z_{i})$ to obtain}{\small\par}
\end{itemize}
\begin{center}
{\small{}$\dfrac{Cov\left(Y_{i},Z_{i}\right)}{Cov\left(S_{i},Z_{i}\right)}=\rho$}{\small\par}
\par\end{center}
\begin{itemize}
\item {\small{}The ratio of the covariance of the outcome with the instrument
to the covariance of the endogenous variable with the instrument identifies
the long regression coefficient $\rho$\medskip{}
}{\small\par}
\item {\small{}This ratio of covariances is the population }\textbf{\small{}instrumental
variables}{\small{} coefficient}{\small\par}
\end{itemize}
\end{frame}

\begin{frame}{Instrumental Variables}

\begin{center}
$\dfrac{Cov\left(Y_{i},Z_{i}\right)}{Cov\left(S_{i},Z_{i}\right)}=\rho$
\par\end{center}
\begin{itemize}
\item Divide numerator and denominator by $Var(Z_{i})$ to obtain
\end{itemize}
\begin{center}
$\dfrac{Cov(Y_{i},Z_{i})/Var(Z_{i})}{Cov(S_{i},Z_{i})/Var(Z_{i})}=\rho$
\par\end{center}
\begin{itemize}
\item What is the interpretation of this expression?
\end{itemize}
\end{frame}

\begin{frame}{Instrumental Variables}

\begin{center}
{\small{}$\dfrac{Cov(Y_{i},Z_{i})/Var(Z_{i})}{Cov(S_{i},Z_{i})/Var(Z_{i})}=\rho$}{\small\par}
\par\end{center}
\begin{itemize}
\item {\small{}The IV coefficient can be viewed as a ratio of two OLS regression
coefficients:}{\small\par}
\end{itemize}
\begin{center}
{\small{}$Y_{i}=\lambda+\delta Z_{i}+u_{i}$}{\small\par}
\par\end{center}

\begin{center}
{\small{}$S_{i}=\kappa+\pi Z_{i}+\eta_{i}$}{\small\par}
\par\end{center}
\begin{itemize}
\item {\small{}$\delta$ is the }\textbf{\small{}reduced form}{\small{}
coefficient, and $\pi$ is the }\textbf{\small{}first stage }{\small{}coefficient\medskip{}
}{\small\par}
\item {\small{}The IV coefficient is $\rho_{IV}=\delta/\pi$, the reduced
form divided by the first stage\medskip{}
}{\small\par}
\item {\small{}If our two IV assumptions are satisfied, then $\rho_{IV}=\rho$:
IV recovers the long regression coefficient}{\small\par}
\end{itemize}
\end{frame}

\begin{frame}{IV Example: Angrist and Krueger (1991)}

\begin{itemize}
\item {\small{}IV eliminates omitted variable bias and allows us to recover
the causal parameter of interest\medskip{}
}{\small\par}
\item {\small{}But where did the instrument $Z_{i}$ come from? \medskip{}
}{\small\par}
\item {\small{}Good instruments typically come from institutional knowledge
combined with plausible assumptions about behavioral relationships\medskip{}
}{\small\par}
\item {\small{}Well-known example: Angrist and Krueger (1991) study of the
returns to education}{\small\par}
\end{itemize}
\end{frame}

\begin{frame}{Angrist and Krueger (1991)}

\begin{itemize}
\item {\small{}AK estimate the causal effect of schooling on earnings\medskip{}
}{\small\par}
\item {\small{}AK's IV strategy is motivated by the interaction between
compulsory schooling and age-at-entry laws\medskip{}
}

\begin{itemize}
\item {\small{}Students can typically drop out of school on the day they
turn 16\medskip{}
}{\small\par}
\item {\small{}Birth date cutoff for starting age: Kids usually start school
in the fall of the calendar year in which they turn six\medskip{}
}{\small\par}
\end{itemize}
\item {\small{}These rules generate differences in mean educational attainment
by date of birth}{\small\par}
\end{itemize}
\end{frame}
\setbeamercolor{background canvas}{bg=} \includepdf{aux/qob_chart.pdf}
\begin{frame}{Angrist and Krueger (1991)}

\begin{itemize}
\item {\small{}The combination of dropout and age-at-entry rules means that
compulsory schooling requirements are higher for those born late in
the year\medskip{}
}{\small\par}
\item {\small{}This suggests an IV strategy using date of birth as an instrument
for schooling\medskip{}
}{\small\par}
\item {\small{}AK use a ``quarter of birth'' instrument, an indicator
for being born early in the year (January - March)\medskip{}
}{\small\par}
\item {\small{}What are the first stage and exclusion restriction assumptions
in this example?}{\small\par}
\end{itemize}
\end{frame}
\setbeamercolor{background canvas}{bg=} \includepdf{aux/angrist_kruger_1991_f1.pdf}

\setbeamercolor{background canvas}{bg=} \includepdf{aux/angrist_krueger_1991_f5.pdf}
\begin{frame}{The Wald Estimator}

\begin{center}
$\rho_{IV}=\dfrac{Cov(Y_{i},Z_{i})/Var(Z_{i})}{Cov(S_{i},Z_{i})/Var(Z_{i})}$
\par\end{center}
\begin{itemize}
\item When $Z_{i}$ is binary, the first stage and reduced form coefficients
become differences in means:
\end{itemize}
\begin{center}
$\rho_{IV}=\dfrac{E\left[Y_{i}|Z_{i}=1\right]-E\left[Y_{i}|Z_{i}=0\right]}{E\left[S_{i}|Z_{i}=1\right]-E\left[S_{i}|Z_{i}=0\right]}$
\par\end{center}
\begin{itemize}
\item This is the population analogue of the Wald (1940) estimator, the
ratio of the mean difference in the outcome to the mean difference
in the endogenous regressor
\end{itemize}
\end{frame}
\setbeamercolor{background canvas}{bg=} \includepdf{aux/angrist_krueger_1991_t3.pdf}

\begin{frame}{IV with Covariates}

\begin{itemize}
\item {\small{}When doing IV it's often necessary to control for additional
observed variables\medskip{}
}{\small\par}
\item {\small{}Example: QOB is correlated with age, which likely also affects
earnings. This suggests a violation of the exclusion restriction;
we may want to control for smooth functions of age\medskip{}
}{\small\par}
\item {\small{}The long regression of interest now includes the control
vector $X_{i}$:}{\small\par}
\end{itemize}
\begin{center}
{\small{}$Y_{i}=\alpha+\rho S_{i}+X_{i}^{\prime}\gamma+\gamma A_{i}+v_{i}$}{\small\par}
\par\end{center}
\begin{itemize}
\item {\small{}Control variables that are neither instruments nor instrumented
are ``exogenous covariates''}{\small\par}
\end{itemize}
\end{frame}

\begin{frame}{IV with Covariates}

\begin{itemize}
\item {\small{}Reduced form and first stage regressions augmented with covariates:}{\small\par}
\end{itemize}
\begin{center}
{\small{}$Y_{i}=\lambda+\delta Z_{i}+X_{i}^{\prime}\psi+u_{i}$}{\small\par}
\par\end{center}

\begin{center}
{\small{}$S_{i}=\kappa+\pi Z_{i}+X_{i}^{\prime}\phi+\eta_{i}$}{\small\par}
\par\end{center}
\begin{itemize}
\item {\small{}As before we can solve for $\rho$ as $\delta/\pi$\medskip{}
}{\small\par}
\item {\small{}Applying the Frisch-Waugh-Lovell theorem, this can be written}{\small\par}
\end{itemize}
\begin{center}
{\small{}$\dfrac{\delta}{\pi}=\dfrac{Cov\left(Y_{i},\tilde{Z}_{i}\right)}{Cov\left(S_{i},\tilde{Z}_{i}\right)}$}{\small\par}
\par\end{center}

\end{frame}

\begin{frame}{Indirect Least Squares}

\begin{center}
{\small{}$\dfrac{\delta}{\pi}=\dfrac{Cov\left(Y_{i},\tilde{Z}_{i}\right)}{Cov\left(S_{i},\tilde{Z}_{i}\right)}$}{\small\par}
\par\end{center}
\begin{itemize}
\item {\small{}The sample analogue of this ratio is called the Indirect
Least Squares (ILS) estimator \medskip{}
}{\small\par}
\item {\small{}Controlling for covariates is equivalent to using the residual
from a regression of $Z_{i}$ on $X_{i}$ as the instrument\medskip{}
}{\small\par}
\item {\small{}The first stage and exclusion restriction assumptions now
need to hold for $\tilde{Z}_{i}$ rather than $Z_{i}$\medskip{}
}{\small\par}
\item {\small{}Importantly, this means that the first stage relationship
between $S_{i}$ and $Z_{i}$ must survive the addition of the controls
$X_{i}$}{\small\par}
\end{itemize}
\end{frame}
\subsection{2SLS and Multiple Instruments}
\begin{frame}{Multiple Instruments }

\begin{itemize}
\item {\small{}Now suppose we have multiple instruments, $Z_{i1}...Z_{iJ}$,
collected in the $J\times1$ vector $Z_{i}$\medskip{}
}{\small\par}
\item {\small{}Example: dummies for each quarter of birth, rather than just
the first quarter\medskip{}
}{\small\par}
\item {\small{}It's no longer clear how to solve for the long regression
coefficient $\rho$. We could use any one of the $J$ instruments,
or any linear combination of them\medskip{}
}{\small\par}
\item {\small{}$\rho$ is now }\textbf{\small{}overidentified}{\small{}:
there are multiple ways to solve for the parameter of interest\medskip{}
}{\small\par}
\item {\small{}How should we combine the available information?}{\small\par}
\end{itemize}
\end{frame}

\begin{frame}{Multiple Instruments }

\begin{itemize}
\item {\small{}Our structural equation of interest is still }{\small\par}
\end{itemize}
\begin{center}
{\small{}$Y_{i}=\alpha+\rho S_{i}+X_{i}^{\prime}\gamma+\epsilon_{i}$}{\small\par}
\par\end{center}
\begin{itemize}
\item {\small{}Consider a first-stage regression of $S_{i}$ on all the
instruments and exogenous controls:}{\small\par}
\end{itemize}
\begin{center}
{\small{}$S_{i}=\kappa+Z_{i}^{\prime}\pi+X_{i}^{\prime}\phi+\eta_{i}$}{\small\par}
\par\end{center}

\end{frame}

\begin{frame}{Multiple Instruments }

\begin{itemize}
\item {\small{}Substitute the first stage regression into the structural
equation:}{\small\par}
\end{itemize}
\begin{center}
{\small{}$Y_{i}=\alpha+\rho\left[\kappa+Z_{i}^{\prime}\pi+X_{i}^{\prime}\phi+\eta_{i}\right]+X_{i}^{\prime}\gamma+\epsilon_{i}$}{\small\par}
\par\end{center}

\begin{center}
{\small{}$=\alpha+\rho\left(\kappa+Z_{i}^{\prime}\pi+X_{i}^{\prime}\phi\right)+X_{i}^{\prime}\gamma+\left(\rho\eta_{i}+\epsilon_{i}\right)$}{\small\par}
\par\end{center}

\begin{center}
{\small{}$=\alpha+\rho\hat{S}_{i}+X_{i}^{\prime}\gamma+\nu_{i}$}{\small\par}
\par\end{center}
\begin{itemize}
\item {\small{}Here $\hat{S}_{i}=\kappa+Z_{i}^{\prime}\pi+X_{i}^{\prime}\phi$
is the predicted value of $S_{i}$ from the first stage\medskip{}
}{\small\par}
\item {\small{}The new residual is $\nu_{i}=\rho\eta_{i}+\epsilon_{i}=\rho(S_{i}-\hat{S}_{i})+\epsilon_{i}$}{\small\par}
\end{itemize}
\end{frame}

\begin{frame}{Two-Stage Least Squares}

\begin{center}
{\small{}$Y_{i}=\alpha+\rho\hat{S}_{i}+X_{i}^{\prime}\gamma+\nu_{i}$}{\small\par}
\par\end{center}
\begin{itemize}
\item {\small{}$\hat{S}_{i}$ and $X_{i}$ are uncorrelated with $\nu_{i}$
(why?), so OLS recovers the parameters of this equation\medskip{}
}{\small\par}
\item {\small{}This is the two-stage least squares (2SLS) estimator:\medskip{}
}

\begin{enumerate}
\item {\small{}Regress $S_{i}$ on all the instruments and exogenous controls,
and obtain predicted values $\hat{S}_{i}$\medskip{}
}{\small\par}
\item {\small{}Regress $Y_{i}$ on $\hat{S}_{i}$ and the exogenous controls\medskip{}
}{\small\par}
\end{enumerate}
\item {\small{}2SLS provides one method for combining instruments. Do you think it is the most efficient?}{\small\par}
\end{itemize}
\end{frame}
\subsection{IV as GMM}
\begin{frame}{IV Estimation}

\begin{itemize}
\item {\small{}Suppose we have $N$ observations on $(Y_{i},S_{i},Z_{i},X_{i})$,
where \medskip{}
}

\begin{itemize}
\item {\small{}$S_{i}$ is a $K\times1$ vector of endogenous variables
(there now may be more than one)\medskip{}
}{\small\par}
\item {\small{}$Z_{i}$ is a $J\times1$ vector of instruments\medskip{}
}{\small\par}
\item {\small{}$X_{i}$ is an $L\times1$ vector of exogenous controls\medskip{}
}{\small\par}
\end{itemize}
\item {\small{}Let $\mathbf{X}_{i}=\left(1,X_{i}^{\prime},S_{i}^{\prime}\right)^{\prime}$
and $\mathbf{Z}_{i}=(1,X_{i}^{\prime},Z_{i}^{\prime})^{\prime}$}{\small\par}
\end{itemize}
\end{frame}


\begin{frame}{IV as GMM}
\begin{itemize}
\item {\small{}Let's augment our structural equation}{\small\par}
\begin{center}
{\small{}$Y_{i}=\alpha+S_{i}^{\prime}\rho+X_{i}^{\prime}\gamma+\epsilon_{i}$}{\small\par}
\par\end{center}
\begin{itemize}
\item {\small{}$S_{i}$ is a $K\times1$ vector of endogenous variables
(there now may be more than one)\medskip{}
}{\small\par}
\item {\small{}$Z_{i}$ is a $J\times1$ vector of instruments ($J\geq K$)\medskip{}
}{\small\par}
\item {\small{}$X_{i}$ is an $L\times1$ vector of exogenous controls\medskip{}
}{\small\par}
\item {\small{}Let $\mathbf{X}_{i}=\left(1,X_{i}^{\prime},S_{i}^{\prime}\right)^{\prime}$
and $\mathbf{Z}_{i}=(1,X_{i}^{\prime},Z_{i}^{\prime})^{\prime}$}{\small\par}
\end{itemize}
\pause
\item We know (from either economic theory or institutional knowledge of $Z_{i}$ and $S_{i}$) that 
\be
\E[\mathbf{Z}_{i}'\varepsilon_{i}]=\E[\mathbf{Z}_{i}'(y_{i}-\alpha-S_{i}^{\prime}\rho-X_{i}^{\prime}\gamma)]=\E[m_{i}(\theta)]=0
\ee
where $\theta=(\rho;\gamma)$
\item What do we call this? How many equations? How many unknowns?
\end{itemize}
\end{frame}
%%
\begin{frame}
\frametitle{Casting everything as GMM}
\begin{itemize}
\item Recall previous class on GMM (see Estimation slide deck). GMM minimizes the following criterion function
\be
\label{GMM_criterion}
Q_{GMM}(\theta)=\E[m_{i}(\theta)'\mathbf{W} m_{i}(\theta)] 
\ee
\item What happens to $\mathbf{W}$ when $J=K$? 
\medskip
\pause
\item How do we choose $\mathbf{W}$ when $J>K$?
\pause
\medskip
\item Turns out that 2SLS minimizes (\ref{GMM_criterion}) by setting $\mathbf{W}=E[\mathbf{Z}_{i}\mathbf{Z}_{i}']^{-1}$ $\rightarrow$ When is 2SLS efficient? {\tiny{See Section 4.2.1 of MHE for full proof.}}
\pause
\medskip
\item What is the optimal $\mathbf{W}$ in a general setup?
\end{itemize}
\end{frame}
%%
%%
\subsection{Overidentification Tests}
\begin{frame}
\frametitle{Over-id Test}
\begin{itemize}
\item GMM formulation also useful to think on how we can test the validity of the ``model". 
\medskip
\item When there are more instruments than endogenous variables ($J>K$) there are multiple ways to solve for the structural parameters. 
\medskip
\item Essentially an over-id specification test (as introduced in the Estimation slide-deck) tells us how close to being satisfied the moment condition are. 
\medskip
\item In the IV setup, this means that an overid test checks whether \textit{all} the IV moment restrictions are close enough to being satisfied. 
\end{itemize}
\end{frame}
%%
\begin{frame}
\frametitle{Overid Tests in the IV world}
\begin{itemize}
\item Let $\hat{\theta}$ be the minimizer of the sample GMM criterion function 
\be
\label{GMM_criterion_sample}\hat{\theta}=\argmin_{c}\hat{Q}_{GMM}(c)=N\hat{m}_{i}(c)'\mathbf{W} \hat{m}_{i}(c) 
\ee
where $\hat{m}_{i}(\theta)=\dfrac{1}{N}\mathbf{Z}_{i}'(y_{i}-\alpha-S_{i}^{\prime}\rho-X_{i}^{\prime}\gamma)$ (evaluated at the ``truth" $\theta$)
\medskip
\item If $\mathbf{W}= \E[\sqrt{N}\hat{m}_{i}(\theta)\sqrt{N}\hat{m}_{i}(\theta)']^{-1}$, then as $N\rightarrow \infty$
\medskip
\be
\hat{Q}_{GMM}(\hat{\theta})\sim\chi^{2}_{J-K}
\ee
(sometimes this is also called the J-test in the IV setup).
\medskip
\item If $\hat{Q}_{GMM}(\hat{\theta})$ is large relative to critical value of $\chi^{2}_{J-K}$ $\rightarrow$ model is rejected. 
\medskip
\item Rejection means that one of the $J$ instrument is false, i.e. there is one instrument for which $\E[Z_{i}\epsilon_{i}]$ is actually not zero!
\end{itemize}
\end{frame}
%%
\begin{frame}
\frametitle{Intuition}

\begin{itemize}
\item Suppose that there are no controls and we have one endogenous variable $S_{i}$ and two potential instruments $Z_{i1}$ and $Z_{i2}$.  
\medskip
\item Two instruments give us two potential (simple) IV estimates of $\rho$. 
\medskip 
\item First estimate of $\rho$ obtained by solving moment condition 
\be
\E[Z_{i1}(y_{i}-S_{i}\rho_{1})]=0
\ee 
\item And the second one is obtained by fitting 
\be
\E[Z_{i2}(y_{i}-S_{i}\rho_{2})]=0
\ee
\pause
\item Over-id tests whether a model that imposes $\rho_{2}=\rho_{1}$ can fit the data. 
\medskip
\item What is the key underlying assumption that we need to make this test truly useful?
\end{itemize}
\end{frame}
\subsection{Imperfect Compliance}
\begin{frame}{Which research questions can be answered causally?}
\begin{itemize}
    \item Go back to our discussion we had in the first class: often times economists are interested in the causal effect of \textit{choices} made by agents (e.g. attend a private university, smoking, do a pre-cancer screening test, etc).
    \medskip
    \item The gold standard in research is the RCT but sometimes you can't randomize individuals into treatment (IRB won't allow it and for very good reasons).
    \medskip
    \item Strategy is then to randomize $Z_{i}$ where $Z_{i}$ is some incentive to take up a given treatment of interest, $D_{i}$.
    \medskip
    \item Even if there is imperfect compliance $\Pr(D_{i}=1| Z_{i}=1)\neq 1$ one can use IV techniques to estimate a causal effect. 
    \medskip
    \item Simple but powerful idea that still needs to make its way in \textit{all} fields...   
\end{itemize}
\end{frame}
\begin{frame}{IV Makes Its Way in Clinical Trials}
\begin{itemize}
    \item RQ: do pre-cancer screenings improve health?
    \medskip
    \item Can't force doctors (or patients) to test so people have tried to answer this question by randomizing the opportunity to screen... 
    \medskip
    \item Studies in this area thus report Intention-to-Treat (or Intention-to-Screen) Effects: $\E[Y_{i}\vert Z_{i}=1]-\E[Y_{i}\vert Z_{i}=0]$.
    \medskip
    \item There seems to be a lot of heterogeneity in the ITTs estimated across a whole variety of studies 
    \medskip
    \item Winawer (2023) asks ``are these intention-to-treat observations applicable to other
clinical environments?"   
\end{itemize}
\end{frame}
\begin{frame}[plain]

\begin{center}
\includegraphics[scale=0.4]{aux/angrist_hull_IV_1}
\par\end{center}


\end{frame}
\begin{frame}[plain]

\begin{center}
\adjustbox{max width=0.95\textwidth,max totalheight=0.85\textheight}{\includegraphics[scale=0.4]{aux/angrist_hull_IV_0}}
\par\end{center}


\end{frame}
\begin{frame}[plain]

\begin{center}
\includegraphics[scale=0.3]{aux/angrist_hull_IV_2}
\par\end{center}


\end{frame}

\subsection{2SLS Mistakes}
\begin{frame}{2SLS Mistakes}

\begin{itemize}
\item {\small{}2SLS is the workhorse IV estimator used by empirical researchers\medskip{}
}{\small\par}
\item {\small{}It's worth covering some common pitfalls related to 2SLS
estimation}{\small\par}
\end{itemize}
\end{frame}

\begin{frame}{2SLS Mistakes 1: Standard Errors}

\begin{itemize}
\item {\small{}It's usually not a good idea to run the two stages of 2SLS
manually\medskip{}
}{\small\par}
\item {\small{}For one thing, this will give the wrong standard errors\medskip{}
}{\small\par}
\item {\small{}Consider a setting with one endogenous variable $S_{i}$}{\small\par}
\end{itemize}
\begin{center}
{\small{}$Y_{i}=\alpha+\rho S_{i}+\epsilon_{i}$}{\small\par}
\par\end{center}
\begin{itemize}
\item {\footnotesize{}If you run 2SLS manually, the second stage is}{\footnotesize\par}
\end{itemize}
\begin{center}
{\footnotesize{}$Y_{i}=\alpha+\rho\hat{S}_{i}+v_{i}$}{\footnotesize\par}
\par\end{center}
\begin{itemize}
\item {\footnotesize{}Here $v_{i}=\epsilon_{i}+\rho(S_{i}-\hat{S}_{i})$\medskip{}
}{\footnotesize\par}
\end{itemize}
\begin{itemize}
\item {\footnotesize{}Manual 2SLS will typically generate standard errors
that are too large.
\medskip{}
}{\footnotesize\par}
\item {\footnotesize{}The residual variance should be calculated from the
structural residual $Y_{i}-\hat{\alpha}-\hat{\rho}S_{i}$, not the
second-stage residual $Y_{i}-\hat{\alpha}-\hat{\rho}\hat{S}_{i}$}{\footnotesize\par}
\end{itemize}
\end{frame}

\begin{frame}{2SLS Mistakes 2: Covariate Ambivalence}

\begin{itemize}
\item {\small{}A second possible mistake in manual 2SLS is the treatment
of covariates\medskip{}
}{\small\par}
\item {\small{}Suppose we run the 2SLS procedure}{\small\par}
\end{itemize}
\begin{center}
{\small{}$S_{i}=\kappa+Z_{i}^{\prime}\pi+X_{1i}^{\prime}\phi_{1}+\eta_{i}$}{\small\par}
\par\end{center}

\begin{center}
{\small{}$Y_{i}=\alpha+\rho\hat{S}_{i}+X_{1i}^{\prime}\theta_{1}+X_{2i}^{\prime}\theta_{2}+v_{i}$}{\small\par}
\par\end{center}
\begin{itemize}
\item {\small{}The covariate vector $X_{2i}$ is included in the second
stage but not the first (the opposite of the instruments $Z_{i}$)\medskip{}
}{\small\par}
\item {\small{}What's wrong with this?}{\small\par}
\end{itemize}
\end{frame}

\begin{frame}{2SLS Mistakes 2: Covariate Ambivalence}

\begin{center}
{\small{}$S_{i}=\kappa+Z_{i}^{\prime}\pi+X_{1i}^{\prime}\phi_{1}+\eta_{i}$}{\small\par}
\par\end{center}

\begin{center}
{\small{}$Y_{i}=\alpha+\rho\hat{S}_{i}+X_{1i}^{\prime}\theta_{1}+X_{2i}^{\prime}\theta_{2}+v_{i}$}{\small\par}
\par\end{center}
\begin{itemize}
\item {\small{}The second-stage residual is $v_{i}=\epsilon_{i}+\rho(S_{i}-\hat{S}_{i})$\medskip{}
}{\small\par}
\item {\small{}Since $S_{i}-\hat{S}_{i}$ is a residual from a regression
that includes $X_{1i},$ $Cov(X_{1i},S_{i}-\hat{S}_{i})=0$\medskip{}
}{\small\par}
\item {\small{}But in general $Cov(S_{i}-\hat{S}_{i},X_{2i})\neq0$, so
we may have created a correlation between the error term and covariates,
leading to bias in all parameter estimates\medskip{}
}{\small\par}
\item {\small{}``If a covariate is good enough for the second stage, it's
good enough for the first.''}{\small\par}
\end{itemize}
\end{frame}

\begin{frame}{2SLS Mistakes 3: The Forbidden Regression}

\begin{itemize}
\item {\small{}In response to repeated mistakes by graduate students, Jerry
Hausman labeled one 2SLS mistake ``the forbidden regression'' \medskip{}
}{\small\par}
\item {\small{}The forbidden regression involves the erroneous application
of 2SLS logic to nonlinear models\medskip{}
}{\small\par}
\item {\small{}Suppose there is a dummy endogenous variable $T_{i}\in\{0,1\}$,
a vector of covariates $X_{i}$ and a vector of instruments $Z_{i}$\medskip{}
}{\small\par}
\item {\small{}The causal equation is}{\small\par}
\end{itemize}
\begin{center}
$Y_{i}=\alpha+\rho T_{i}+X_{i}^{\prime}\theta+\epsilon_{i}$
\par\end{center}

\end{frame}

\begin{frame}{2SLS Mistakes 3: The Forbidden Regression}

\begin{itemize}
\item {\small{}First stage regression:}{\small\par}
\end{itemize}
\begin{center}
{\small{}$T_{i}=\kappa+Z_{i}^{\prime}\pi+X_{i}^{\prime}\phi+\eta_{i}$}{\small\par}
\par\end{center}
\begin{itemize}
\item {\small{}Since $T_{i}$ is binary, $E\left[T_{i}|Z_{i},X_{i}\right]$
is likely nonlinear unless the model is saturated. The first stage
is therefore a misspecified model of the CEF \medskip{}
}{\small\par}
\item {\small{}One might be tempted to try a nonlinear model to better approximate
the CEF, e.g. a probit:}{\small\par}
\end{itemize}
\begin{center}
{\small{}$Pr\left[T_{i}=1|X_{i},Z_{i}\right]=\Phi\left(\kappa+Z_{i}^{\prime}\pi+X_{i}^{\prime}\phi\right)$}{\small\par}
\par\end{center}

\end{frame}

\begin{frame}{2SLS Mistakes 3: The Forbidden Regression}

\begin{itemize}
\item {\small{}The nonlinear first stage can be used to obtain predicted
values:}{\small\par}
\end{itemize}
\begin{center}
{\small{}$\hat{T}_{i}=\Phi\left(\hat{\kappa}+Z_{i}^{\prime}\hat{\pi}+X_{i}^{\prime}\hat{\phi}\right)$}{\small\par}
\par\end{center}
\begin{itemize}
\item {\small{}One might then be tempted to run a 2SLS-style second stage:}{\small\par}
\end{itemize}
\begin{center}
{\small{}$Y_{i}=\alpha+\rho\hat{T}_{i}+X_{i}^{\prime}\theta+\epsilon_{i}$$ $}{\small\par}
\par\end{center}
\begin{itemize}
\item {\small{}This regression is forbidden!}{\small\par}
\end{itemize}
\end{frame}

\begin{frame}{2SLS Mistakes 3: The Forbidden Regression}

\begin{itemize}
\item {\footnotesize{}To see why the forbidden regression is forbidden,
rewrite the structural equation}{\footnotesize\par}
\end{itemize}
\begin{center}
{\footnotesize{}$Y_{i}=\alpha+\rho\hat{T}_{i}+X_{i}^{\prime}\theta+\left(\epsilon_{i}+\rho(T_{i}-\hat{T}_{i})\right)$}{\footnotesize\par}
\par\end{center}
\begin{itemize}
\item {\footnotesize{}If the first stage is estimated via OLS, then the
residual $(T_{i}-\hat{T}_{i})$ is guaranteed to be uncorrelated with
the fit $\hat{T}_{i}$ and with the included covariates $X_{i}$\medskip{}
}{\footnotesize\par}
\item {\footnotesize{}Nonlinear residuals do not have this property. $Cov\left(T_{i}-\hat{T}_{i},\hat{T}_{i}\right)$
and $Cov\left(T_{i}-\hat{T}_{i},X_{i}\right)$ will generally only
equal zero when the $E\left[T_{i}|X_{i},Z_{i}\right]$ is really a
probit\medskip{}
}{\footnotesize\par}
\item {\footnotesize{}Who knows whether the probit CEF is right? Better
to stick with a linear first stage, which guarantees zero covariance
between residuals and fits/regressors even if the CEF is misspecified}{\footnotesize\par}
\end{itemize}
\end{frame}
\begin{frame}
\frametitle{Pop-Quiz}
\pause
\begin{itemize}
\item Suppose you came up with what you think is a good instrument, $Z_{i}$.
\medskip
\item You are however not 100\% sure about the exclusion restriction.
\medskip
\item Inspired by your knowledge of regression and the meaning of the exclusion restriction you come up with the following test for the validity of the instrument
\be
y_{i}=\alpha + \beta S_{i} + \gamma Z_{i} + r_{i}
\ee
\item That is you run a regression on your endogenous variable $S_{i}$ and the instrument $Z_{i}$. 
\medskip
\item Since the instrument only operates through the endogenous variable $S_{i}$, you argue that if the exclusion restriction is valid, then $\beta=\rho$ (the true casual effect) and $\gamma=0$.
\medskip
\item Right? Wrong?
\end{itemize}
\end{frame}

\begin{frame}
\centering
\LARGE{APPENDIX}
\end{frame}
%%%%%%%%%%%%%%%%%%%%%
%%%%%%%%%%%%%%%%%%%%%

\begin{frame}{IV in Finite Samples}

\begin{itemize}
\item {\small{}OLS is an unbiased
estimator of population projection coefficients in finite samples
(with a linear CEF or nonstochastic regressors)\medskip{}
}{\small\par}
\item {\small{}IV estimators do not have this property\medskip{}
}{\small\par}
\item {\small{}IV is a ratio of reduced form and first stage coefficients:}{\small\par}
\end{itemize}
\begin{center}
{\small{}$E\left[\hat{\rho}_{IV}\right]=E\left[\dfrac{\hat{\delta}}{\hat{\pi}}\right]$}{\small\par}
\par\end{center}

\begin{center}
{\small{}$\neq\dfrac{E\left[\hat{\delta}\right]}{E\left[\hat{\pi}\right]}$}{\small\par}
\par\end{center}

\begin{center}
{\small{}$=\rho$}{\small\par}
\par\end{center}

\end{frame}

\begin{frame}{IV in Finite Samples: Example}

\begin{itemize}
\item {\footnotesize{}Consider the following reduced form and first stage:}{\footnotesize\par}
\end{itemize}
\begin{center}
{\footnotesize{}$Y_{i}=\lambda+\delta Z_{i}+u_{i}$}{\footnotesize\par}
\par\end{center}

\begin{center}
{\footnotesize{}$S_{i}=\kappa+\pi Z_{i}+\eta_{i}$}{\footnotesize\par}
\par\end{center}
\begin{itemize}
\item {\footnotesize{}Suppose the instruments are non-stochastic, and }{\footnotesize\par}
\end{itemize}
\begin{center}
{\footnotesize{}$(u_{i},\eta_{i})\sim N\left(0,\left[\begin{array}{cc}
\sigma_{u}^{2} & \tau\sigma_{u}\sigma_{\eta}\\
\tau\sigma_{u}\sigma_{\eta} & \sigma_{\eta}^{2}
\end{array}\right]\right)$}{\footnotesize\par}
\par\end{center}
\begin{itemize}
\item {\footnotesize{}Then the finite-sample distribution of the reduced
form and first stage coefficients is}{\footnotesize\par}
\end{itemize}
\begin{center}
{\footnotesize{}$\left(\hat{\delta},\hat{\pi}\right)\sim N\left((\delta,\pi),\left[\begin{array}{cc}
\sigma_{u}^{2}\left[\sum_{i}(Z_{i}-\bar{Z})^{2}\right]^{-1} & \tau\sigma_{u}\sigma_{\eta}\left[\sum_{i}(Z_{i}-\bar{Z})^{2}\right]^{-1}\\
\tau\sigma_{u}\sigma_{\eta}\left[\sum_{i}(Z_{i}-\bar{Z})^{2}\right]^{-1} & \sigma_{\eta}^{2}\left[\sum_{i}(Z_{i}-\bar{Z})^{2}\right]^{-1}
\end{array}\right]\right)$}{\footnotesize\par}
\par\end{center}

\end{frame}

\begin{frame}{IV in Finite Samples: Example}

\begin{itemize}
\item The IV coefficient, $\hat{\rho}=\hat{\delta}/\hat{\pi}$, is a ratio
of two jointly normal random variables{\small{}\medskip{}
}{\small\par}
\item $\hat{\rho}$ therefore follows a normal ratio distribution {\small{}\medskip{}
}{\small\par}
\item Normal ratio distributions do not have moments, so $E\left[\hat{\rho}\right]$,
$Var\left(\hat{\rho}\right)$, and all higher moments do not exist.
In this sense $\hat{\rho}$ is biased{\small{}\medskip{}
}{\small\par}
\item If $\tau=0$ and $\pi\rightarrow0$, $\hat{\rho}$ becomes the ratio
of two independent mean zero normals, which follows a Cauchy distribution
(see Nelson and Starz 1990 for more)
\end{itemize}
\end{frame}

\begin{frame}{IV in Finite Samples: Many/Weak Problems}

\begin{itemize}
\item {\footnotesize{}In general, just-identified IV estimators have no
moments and fat tails\medskip{}
}{\footnotesize\par}
\item {\footnotesize{}However, just-id estimators are approximately centered
in the right place in the sense that $Q_{0.5}(\hat{\rho}_{IV})\approx\rho$
(more details later)\medskip{}
}{\footnotesize\par}
\item {\footnotesize{}It turns out that overidentified IV estimators are
not centered at the right place: they are biased towards OLS\medskip{}
}{\footnotesize\par}
\item {\footnotesize{}The severity of finite-sample bias increases with
the degree of overfitting AND as the first-stage relationship
between instruments and endogenous regressors gets weaker\medskip{}
}{\footnotesize\par}
\item {\footnotesize{}IV bias is therefore worse when the instruments are
``many/weak''\medskip{}
}{\footnotesize\par}
\item {\footnotesize{}Bound, Jaeger and Baker (1995) showed that this problem
can be important in practice, spawning a large literature on dealing
with many/weak instruments}{\footnotesize\par}
\end{itemize}
\end{frame}

\begin{frame}{Bound, Jaeger and Baker (1995)}

\begin{itemize}
\item {\footnotesize{}Bound, Jaeger and Baker (1995) revisit the Angrist/Krueger
QOB instrument\medskip{}
}{\footnotesize\par}
\item {\footnotesize{}The Wald estimates we studied are just-identified,
so the many/weak critique does not apply (a common misconception)
\medskip{}
}{\footnotesize\par}
\item {\footnotesize{}However, AK also estimate overidentified 2SLS models
of the form}{\footnotesize\par}
\end{itemize}
\begin{center}
{\footnotesize{}$Y_{i}=\rho S_{i}+\sum_{b=1}^{50}\theta_{b}B_{ib}+\sum_{t=1}^{30}\psi_{t}T_{it}+\epsilon_{i}$}{\footnotesize\par}
\par\end{center}

\begin{center}
{\footnotesize{}$S_{i}=\sum_{q=1}^{3}\left[\sum_{b=1}^{50}\pi_{qb}B_{ib}Q_{iq}+\sum_{t=1}^{30}\tau_{qt}T_{it}Q_{iq}\right]$}{\footnotesize\par}
\par\end{center}

\begin{center}
{\footnotesize{}$+\sum_{b=1}^{50}\kappa_{b}B_{ib}+\sum_{t=1}^{30}\phi_{t}T_{it}+\eta_{i}$}{\footnotesize\par}
\par\end{center}
\begin{itemize}
\item {\footnotesize{}Here $B_{ib}$ is a dummy for birth in state $b$,
$T_{it}$ is a dummy for birth in year $t$, and $Q_{iq}$ is a dummy
for birth in quarter $q$}{\footnotesize\par}
\end{itemize}
\end{frame}

\begin{frame}{Bound, Jaeger and Baker (1995)}

\begin{itemize}
\item {\footnotesize{}The AK model uses 30 quarter/year and 150 quarter/state
interactions as instruments, so is overidentified by 179\medskip{}
}{\footnotesize\par}
\item {\footnotesize{}The interactions are meant to capture differences
in the effects of compulsory schooling laws across states and over
time\medskip{}
}{\footnotesize\par}
\item {\footnotesize{}From the point of view of 2SLS asymptotics, adding
the interactions can only be a plus: more instruments improve the
first stage fit and drive down standard errors\medskip{}
}{\footnotesize\par}
\item {\footnotesize{}But the QOB first stage is small and many of the interactions
are likely to be irrelevant, so this is a severe many/weak scenario\medskip{}
}{\footnotesize\par}
\item {\footnotesize{}Bound, Jaeger and Baker replace the QOB dummies with
randomly-generated fake QOB variables that have no first stage by
design}{\footnotesize\par}
\end{itemize}
\end{frame}

\begin{frame}{Bound, Jaeger and Baker (1995)}

\begin{center}
{\footnotesize{}}%
\begin{tabular}{|c|c|c|c|}
\hline 
 & {\footnotesize{}2SLS (real QOB)} & {\footnotesize{}2SLS (fake QOB)} & {\footnotesize{}OLS}\tabularnewline
\hline 
\hline 
{\footnotesize{}30 QOB{*}year} & {\footnotesize{}0.081} & {\footnotesize{}0.062} & {\footnotesize{}0.063}\tabularnewline
\hline 
{\footnotesize{}interactions} & {\footnotesize{}(0.016)} & {\footnotesize{}(0.037)} & {\footnotesize{}(0.000)}\tabularnewline
\hline 
 &  &  & \tabularnewline
\hline 
{\footnotesize{}180 QOB{*}year and state} & {\footnotesize{}0.083} & {\footnotesize{}0.060} & {\footnotesize{}0.063}\tabularnewline
\hline 
{\footnotesize{}interactions} & {\footnotesize{}(0.009)} & {\footnotesize{}(0.015)} & {\footnotesize{}(0.000)}\tabularnewline
\hline 
\end{tabular}{\footnotesize\par}
\par\end{center}

\end{frame}

\begin{frame}{Bound, Jaeger and Baker (1995)}

\begin{itemize}
\item {\small{}The Bound, Jaeger and Baker results show that the AK 2SLS
specification essentially reproduces OLS even if the instruments are
junk unrelated to schooling\medskip{}
}{\small\par}
\item {\small{}Disturbingly, conventional IV standard errors give no indication
of the extent of the problem\medskip{}
}{\small\par}
\item {\small{}The fake QOB 180-instrument model generates an estimate that
looks precise despite the lack of a first stage\medskip{}
}{\small\par}
\item {\small{}How can we diagnose many/weak problems, and correct them
if they are present?}{\small\par}
\end{itemize}
\end{frame}

\begin{frame}{Many/Weak Bias}

\begin{itemize}
\item {\small{}Consider a model with a single endogenous regressor and $Q$
instruments, written in matrix form:}{\small\par}
\end{itemize}
\begin{center}
$Y=\rho S+\epsilon$
\par\end{center}

\begin{center}
$S=Z\pi+\eta$
\par\end{center}
\begin{itemize}
\item Remember that we can write $\hat{S}=P_{Z}S$
\item We can thus re-write the 2SLS estimator of $\rho$ (regression of $Y$ on $\hat{S}$) as 
\end{itemize}
\begin{center}
$\hat{\rho}_{2SLS}=\rho+(S^{\prime}P_{Z}S)^{-1}S^{\prime}P_{Z}\epsilon$
\par\end{center}

\begin{center}
$=\rho+\left(S^{\prime}P_{Z}S\right)^{-1}(\pi^{\prime}Z^{\prime}+\eta^{\prime})P_{Z}\epsilon$
\par\end{center}

\begin{center}
$=\rho+(S^{\prime}P_{Z}S)^{-1}\pi^{\prime}Z^{\prime}\epsilon+\left(S^{\prime}P_{Z}S\right)^{-1}\eta^{\prime}P_{Z}\epsilon$
\par\end{center}

\end{frame}

\begin{frame}{Many/Weak Bias}

\begin{center}
{\small{}$E\left[\hat{\rho}_{2SLS}-\rho\right]=E\left[(S^{\prime}P_{Z}S)^{-1}\pi^{\prime}Z^{\prime}\epsilon\right]+E\left[\left(S^{\prime}P_{Z}S\right)^{-1}\eta^{\prime}P_{Z}\epsilon\right]$}{\small\par}
\par\end{center}
\begin{itemize}
\item {\small{}It's difficult to make further progress since we cannot pass
expectations inside the inverse\medskip{}
}{\small\par}
\item {\small{}Bekker (1994): passing the expectations through provides
an accurate approximation under group asymptotics that allow $Q$
to grow with $N$\medskip{}
}{\small\par}
\item {\small{}Then the finite-sample bias of 2SLS can be approximated}{\small\par}
\end{itemize}
\begin{center}
{\small{}$E\left[\hat{\rho}_{2SLS}-\rho\right]\approx(E\left[S^{\prime}P_{Z}S\right])^{-1}\pi^{\prime}E\left[Z^{\prime}\epsilon\right]+\left(E\left[S^{\prime}P_{Z}S\right]\right)^{-1}E\left[\eta^{\prime}P_{Z}\epsilon\right]$}{\small\par}
\par\end{center}

\end{frame}

\begin{frame}{Many/Weak Bias}

\begin{center}
{\small{}$E\left[\hat{\rho}_{2SLS}-\rho\right]\approx(E\left[S^{\prime}P_{Z}S\right])^{-1}\pi^{\prime}E\left[Z^{\prime}\epsilon\right]+\left(E\left[S^{\prime}P_{Z}S\right]\right)^{-1}E\left[\eta^{\prime}P_{Z}\epsilon\right]$}{\small\par}
\par\end{center}
\begin{itemize}
\item {\small{}The first term is zero (why)?\medskip{}
}{\small\par}
\item {\small{}The finite-sample bias of 2SLS comes from the fact that the
numerator of the second term is not zero:\medskip{}
}{\small\par}
\end{itemize}
\begin{center}
$E\left[\eta^{\prime}P_{Z}\epsilon\right]\neq0$
\par\end{center}
\begin{itemize}
\item What is the intuition for this?
\end{itemize}
\end{frame}


\begin{frame}{Many/Weak Bias}

\begin{itemize}
\item {\footnotesize{}To make further progress, let's use some matrix tricks}

\begin{itemize}
\item {\footnotesize{}A scalar equals its trace}{\footnotesize\par}
\item {\footnotesize{}Trace is invariant to cyclic permutations}{\footnotesize\par}
\item {\footnotesize{}Trace is a linear operator that passes through expectations}{\footnotesize\par}
\end{itemize}
\end{itemize}
\begin{center}
{\footnotesize{}$E\left[\eta^{\prime}P_{Z}\epsilon|Z\right]=E\left[tr\left(\eta^{\prime}P_{Z}\epsilon\right)|Z\right]$}{\footnotesize\par}
\par\end{center}

\begin{center}
{\footnotesize{}$=E\left[tr\left(P_{Z}\epsilon\eta^{\prime}\right)|Z\right]$}{\footnotesize\par}
\par\end{center}

\begin{center}
{\footnotesize{}$=tr\left(P_{Z}E\left[\epsilon\eta^{\prime}|Z\right]\right)$}{\footnotesize\par}
\par\end{center}

\begin{center}
{\footnotesize{}$=\sigma_{\eta\epsilon}tr\left(Z(Z^{\prime}Z)^{-1}Z^{\prime}\right)$}{\footnotesize\par}
\par\end{center}

\begin{center}
{\footnotesize{}$=\sigma_{\eta\epsilon}tr\left((Z^{\prime}Z)^{-1}Z^{\prime}Z\right)$}{\footnotesize\par}
\par\end{center}

\begin{center}
{\footnotesize{}$=\sigma_{\eta\epsilon}Q$}{\footnotesize\par}
\par\end{center}

\end{frame}

\begin{frame}{Many/Weak Bias}

\begin{itemize}
\item {\footnotesize{}The denominator of the bias is}{\footnotesize\par}
\end{itemize}
\begin{center}
{\footnotesize{}$E\left[S^{\prime}P_{Z}S\right]=E\left[(\pi^{\prime}Z^{\prime}+\eta^{\prime})P_{Z}(Z\pi+\eta)\right]$}{\footnotesize\par}
\par\end{center}

\begin{center}
{\footnotesize{}$=E\left[\pi^{\prime}Z^{\prime}Z\pi\right]+E\left[\eta^{\prime}P_{Z}\eta\right]$}{\footnotesize\par}
\par\end{center}
\begin{itemize}
\item {\footnotesize{}Applying the same trace tricks to the second term
shows that $E\left[\eta^{\prime}P_{Z}\eta\right]=\sigma_{\eta}^{2}Q$}{\small{}\medskip{}
}{\small\par}
\item {\footnotesize{}Then}{\footnotesize\par}
\end{itemize}
\begin{center}
{\footnotesize{}$E\left[\hat{\rho}_{2SLS}-\rho\right]\approx\dfrac{\sigma_{\eta\epsilon}Q}{E\left[\pi^{\prime}Z^{\prime}Z\pi\right]+\sigma_{\eta}^{2}Q}$}{\footnotesize\par}
\par\end{center}

\begin{center}
{\footnotesize{}$=\dfrac{\sigma_{\eta\epsilon}}{\sigma_{\eta^{2}}}\cdot\dfrac{1}{1+E\left[\pi^{\prime}Z^{\prime}Z\pi\right]/(\sigma_{\eta}^{2}Q)}$}{\footnotesize\par}
\par\end{center}

\end{frame}

\begin{frame}{Many/Weak Bias}

\begin{center}
{\small{}$E\left[\hat{\rho}_{2SLS}-\rho\right]\approx\dfrac{\sigma_{\eta\epsilon}}{\sigma_{\eta^{2}}}\cdot\dfrac{1}{1+E\left[\pi^{\prime}Z^{\prime}Z\pi\right]/(\sigma_{\eta}^{2}Q)}$}{\small\par}
\par\end{center}
\begin{itemize}
\item {\small{}Notice that $E\left[\pi^{\prime}Z^{\prime}Z\pi\right]$ is
the explained sum of squares from the regression of $S$ on $Z$,
while $\sigma_{\eta}^{2}$ is the residual sum of squares divided
by $N$\medskip{}
}{\small\par}
\item {\small{}Then }{\small\par}
\end{itemize}
\begin{center}
{\small{}$E\left[\hat{\rho}_{2SLS}-\rho\right]\approx\dfrac{\sigma_{\eta\epsilon}}{\sigma_{\eta^{2}}}\cdot\dfrac{1}{1+F}$}{\small\par}
\par\end{center}
\begin{itemize}
\item {\small{}Here $F$ is the first stage (population) $F$-statistic for the joint
significance of all the instruments}{\small\par}
\end{itemize}
\end{frame}

\begin{frame}{Many/Weak Bias}

\begin{center}
{\small{}$E\left[\hat{\rho}_{2SLS}-\rho\right]=\dfrac{\sigma_{\eta\epsilon}}{\sigma_{\eta^{2}}}\cdot\dfrac{1}{1+F}$}{\small\par}
\par\end{center}
\begin{itemize}
\item {\small{}Note also that when $F$ goes to zero, the 2SLS bias becomes
$\sigma_{\eta\epsilon}/\sigma_{\eta}^{2}$\medskip{}
}{\small\par}
\item {\small{}With a zero first stage, $S=\eta$ so $\sigma_{\eta\epsilon}/\sigma_{\eta}^{2}=Cov(S_{i},\epsilon_{i})/Var(S_{i})$,
the bias of OLS\medskip{}
}{\small\par}
\item {\small{}As the first stage shrinks to zero, all of the variance in
the fitted values comes from the endogenous part of $\hat{S}$, and
the bias of 2SLS equals the bias of OLS}{\small\par}
\end{itemize}
\end{frame}

\begin{frame}{Many/Weak In Practice}

\begin{itemize}
\item {\small{}In practice it is a good idea to look at the first stage
$F$-statistic \medskip{}
}{\small\par}
\item {\small{}In models with exogenous controls this should be the }\emph{\small{}partial
$F$ }{\small{}on the excluded instruments\medskip{}
}{\small\par}
\item {\small{}Stock and Staiger (1997) suggest that many weak typically
becomes a problem when $F<10$, though this is not a theorem\medskip{}
}{\small\par}
\item {\small{}Standard errors based on Bekker (1994) asymptotics are more
accurate in many/weak scenarios (see also Hansen et al. 2008)}{\small\par}
\end{itemize}
\end{frame}

\begin{frame}{Many/Weak In Practice}

\begin{itemize}
\item {\footnotesize{}The Limited Information Maximum Likelihood (LIML)
estimator coincides with 2SLS for just-id models and is approximately
median unbiased for overidentified models \medskip{}
}{\footnotesize\par}
\item {\footnotesize{}A comparison of LIML and 2SLS can help determine whether
many/weak is a problem (though conventional LIML standard errors are
too small) \medskip{}
}{\footnotesize\par}
\item {\footnotesize{}LIML has some unattractive features, however: it has
fat tails (no moments), and is inconsistent with heteroskedasticity\medskip{}
}{\footnotesize\par}
\item {\footnotesize{}There are other estimators (e.g. Fuller) that correct
some of the deficiencies of LIML; see Stock and Yogo (2005)\medskip{}
}{\footnotesize\par}
\item {\footnotesize{}The Anderson-Rubin (AR) test provides a method for
testing hypotheses about $\rho$ that is robust to weak instruments,
though it has low power when the instruments are weak}{\footnotesize\par}
\end{itemize}
\end{frame}
%%
\begin{frame}{The Judges Design}

\begin{itemize}
\item {\footnotesize{}But when it sensible to talk about a research design with ``many" instruments?
$\rightarrow$ a research design that uses differences
in the harshness/leniency of judges or other examiners is a research design that often has to deal with many/weak potential problems\medskip{}
}{\footnotesize\par}
\item {\footnotesize{}Examples: Kling (2006) on criminal sentencing, Doyle
(2007) on foster care, Maestas et al. (2013) on disability insurance,
Dobbie (2015) on bankruptcy \medskip{}
}{\footnotesize\par}
\item {\footnotesize{}If judges are randomly assigned, the identity of the
judge is a plausible instrument (there are many fixed effects for judges!)\medskip{}
}{\footnotesize\par}
\end{itemize}
\end{frame}

\begin{frame}{The Judges Design}

\begin{itemize}
\item {\small{}Common approach to the judges design: Define}{\small\par}
\end{itemize}
\begin{center}
{\small{}$T_{i}^{*}=\dfrac{1}{n_{j(i)}-1}{\displaystyle \sum_{k\neq i}1\left\{ j(k)=j(i)\right\} T_{k}}$}{\small\par}
\par\end{center}
\begin{itemize}
\item {\small{}$T_{i}^{*}$ is the mean treatment rate (e.g. incarceration
rate) for the judge assigned to $i$, $j(i)$, for all cases not involving
$i$\medskip{}
}{\small\par}
\item {\small{}$T_{i}^{*}$ is then used as an instrument for $i$'s actual
treatment $T_{i}$}{\small\par}
\end{itemize}
\end{frame}

\begin{frame}{The Judges Design: JIVE}

\begin{itemize}
\item {\small{}Note that $T_{i}^{*}$ is the fitted value from a first stage
regression of removal on all the judge dummies, leaving out $i$\medskip{}
}{\small\par}
\item {\small{}The instruments are implicitly a full set of judge dummies,
which may lead to many/weak problems\medskip{}
}{\small\par}
\item {\small{}Using a leave-out mean changes the estimator from 2SLS to
jackknife IV (JIVE), which is one solution to the many/weak problem
but not necessarily the optimal one\medskip{}
}{\small\par}
\item {\small{}Conventional standard errors for JIVE are also incorrect
in many/weak scenarios.}{\small\par}
\medskip
\item {\small{}Despite this, many papers use judges design given its implicit attractiveness plus the possibility to summarize nicely heterogenous treatment effects (see the marginal treatment effects literature pioneered by Heckman and Vytacil, 1999,2005).  }{\small\par}
\end{itemize}
\end{frame}



\end{document}